\section{Einleitung}\label{sec:1}

Die vorliegende Arbeit zeigt, wie sich verbale Semantik in Verfassungstexten von 1848 bis 1968 verändert hat. So soll aus rechtslinguistischer Perspektive Evidenz dafür gefunden werden, dass Sprachwandel in Verfassungen mit einem Paradigmenwechsel im Verfassungsrecht korreliert. Die zu beantwortende Forschungsfrage lautet demnach:

\begin{quote}
Ist es möglich, verfassungsrechtliche Paradigmenwechsel durch eine systemisch-funktionale Analyse von Verbphrasen in Verfassungstexten zu identifizieren? Falls ja, wie lassen sich diese grammatisch klassifizieren?
\end{quote}


Somit wird das Thema in das interdisziplinäre Feld der Rechtslinguistik verortet. Dieses ist gerade im Rahmen semantischer Untersuchungen ein wenig erforschter Bereich. Mit dem \textit{konkorp}-Korpus \citep{Thurau2025} liegt nun die Datengrundlage für linguistische Untersuchungen an deutschen Verfassungstexten vor. Das Korpus enthält insgesamt sechs Verfassungstexte aus den Jahren 1848, 1871, 1919, 1949 und 1968. Diese wurden tokenisiert, normalisiert und auf morphosyntaktischer, rechtwissenschaftlich-struktureller und semantischer Ebenen annotiert.

Die Annotationen im \textit{konkorp} ermöglichen es, die zugrundeliegenden Verfassungstexte diachron nach Sprachwandelprozessen zu untersuchen. Besonders von Interesse für die gepalnte Arbeit sind Verbphrasen. Jedoch sollen diese nicht auf formal-syntaktischer Ebene untersucht werden, sondern im Rahmen der \textit{Systemic Functional Linguistics} (SFL, siehe \citealt{Halliday1978,Halliday2004}).

Halliday entwickelte mit der SFL ein Framework, das -- ähnlich wie die Konstruktionsgrammatik (CxG) -- Sprache als ein gebrauchsbasiertes System versteht, welches Bedeutung durch Kontext und Funktionalität darstellt \citep{FillmoreKay1995}.

SFL basiert im Kontrast zur CxG allerdings weniger auf kognitiven als auf soziolinguistischen Konzepten und stellt Sprache als sozial-funktionales System dar. Gerade für die Untersuchung der verbalen Semantik nutzt die SFL Konzepte zur Satzanalyse, die für die Forschung an Verfassungstexten gut geeignet zu sein scheinen. Dieser Transfer ist Hauptgegenstand der geplanten Arbeit. Das System, dass der Verbphrasenanalyse in der SFL zugrundeliegt, ist das \textit{transitivity system} \citep[170]{Halliday2004}. Anders als in formal-syntaktischen Analysen wird Transitivität hier nicht als eine mögliche Verbvalenz verstanden, sondern als semantisches System. Sätze (in der SFL \textit{clauses}) haben drei funktionale Komponenten; den Vorgang (\textit{process}), die am Vorgang Beteiligten (\textit{participants}) und die Umstände des Vorganges (\textit{circumstances}). Der Vorgang wird durch das Verb im Satz ausgedrückt und kann einer von sechs Kategorien (material, relational, behavioural, metal, verbal oder existential) zugeordnet werden \citep[171]{Halliday2004}.

Das Framework eignet sich - nicht zuletzt aufgrund dieser Kategorisierung -, um es auf schriftsprachliche Daten anzuwenden. Gleichzeitig ist die interdisziplinäre Forschung von Verfassungstexten, gerade im Bezug auf Verbsemantik, noch nicht hinreichend erforscht. \cite{Busse2015,BusseEtAl2018} führten Frame-semantische Studien an Gesetzestexten durch und untersuchten dabei Nomina aus dem Zivil- und Strafrecht, \cite{Cirko2022} forschte an Nomina in Grundrechten des deutschen Grundgesetzes und der polnischen Verfassung, und \cite{Sayatz1991} schrieb einen Artikel zur illokutiven Interpretation von Deklarativsätzen mit Modalverben - jedoch findet sich keine korpusbasierte Studie zur diachroner Rechtslinguistik in Bezug auf Verben in Verfassungstexten. Dabei ist die Relevanz grammatischer Analysen für die Rechtswissenschaft nicht abzustreiten. Seit den 1980ern untersucht die sogennante Heidelberger Gruppe rechtslinguistische Themen, und auch international ist die Disziplin der \textit{Legal linguistics} in der Sprachwissenschaft gut verankert. Demnach soll die geplante Arbeit auch dazu dienen, die oben beschriebene Forschungslücke zu füllen. Durch einen systemisch-funktionalen Ansatz soll ein Framework genutzt werden, das bis dato für die Rechtslinguistik noch nicht verwendet worden ist. Sollte es gelingen, SFL auf die Rechtslinguistik zu übertragen, kann dies ein Schritt Richtung methodologischer Vielfalt zur Erschließung der Semantik von Gesetzestexten sein.

Der Aufbau der Arbeit ist aufgebaut wie folgt:

In Abschnitt \ref{sec:2} werden die zugrundeliegenden theoretischen Konzepte eingeführt. Zunächst wird in Abschnitt \ref{sec:2:sfl} die \textit{Systemic Functional Linguistics} (SFL) nach Michael Halliday vorgestellt. In diesem Rahmen werden die Grundzüge der Verbanalyse in der SFL charakterisiert und die vier Ebenen (\textit{Phonology}, \textit{Context}, \textit{Lexico-grammar} und \textit{Semantics}), auf denen sprachliche Äußerungen (auch in Textform) analysiert werden können, erklärt. Bis auf die Phonologie werden in die geplante Analyse alle Ebenen mit einbezogen. Ebenso wird kurz erläutert, warum die Datengrundlage Anreize zum Arbeiten in diesem Framework schafft.

Besonderer Fokus wird auf das Konzept der Transitivität gelegt. Anders als in der formalen Syntak wird Transitivität in der SFL als eine funktionale Kategorie zur Erschließung von Satzbedeutungen betrachtet. Diese setzt sich zusammen aus a) den Vorgängen, die beschrieben werden, b) den an den Vorgängen Beteiligten und c) den Umständen der Vorgänge. Später wird in der Methodik beschrieben, wie dieses Konzept zur qualitativen Analyse von Verbclustern in der vorliegenden Datengrundlage genutzt wird.

In Abschnitt \ref{sec:2:ll} wird ein Überblick über die Forschungsschwerpunkte der Rechtslinguistik gegeben. Anders als in \citealt{Thurau2025} ist geplant, Rechtslinguistik auch kurz von der forensischen Linguistik abzugrenzen. %Jedoch möchte ich, ähnlich wie \citeauthor{Thurau2025}, durch die Arbeit betonen, dass Rechtssprache ein eigenes Register skizziert.

Für den nötigen interdisziplinären Kontext folgt in Abschnitt \ref{sec:2:rg} eine kurze Darstellung der Rechtsgeschichte \citep{Hähnchen2016}. Um das rechtshistorische Fundament für die Datengrundlage der Studie zu schaffen, wird in Abschnitt \ref{sec:2:vr} ein Fokus auf das deutsche Verfassungsrecht gelegt \citep{WeberFas1983,IpsenEtAl2022}.

Anschließend wird in Abschnitt \ref{sec:2:kl} die Disziplin der Korpuslinguistik anhand der Datengrundlage für diese Arbeit vorgestellt und die Vorteile von multi-layered Korpora erläutert \citep{OdebrechtEtAl2016}. Das zu untersuchende Korpus ist hierbei das \textit{konkorp}. Speziell werden Studien aus \citealt{VogelEtAl2022} genutzt, um den Nutzen interdisziplinärer Forschung zwischen der Jurisprudenz und Korpuslinguistik darzustellen.

%Ziel ist es, die genannten Konzepte in der qualitativen Analyse so miteinander zu verknüpfen, dass die oben genannte Forschungsfrage beantwortet werden kann. Im Fokus wird dabei der methodologische Versuch des Transfers der SFL auf die Korpuslinguistik am Beispiel des \textit{konkorp} sein.

In Abschnitt \ref{sec:3:konkorp} wird das \textit{konkorp} detaillierter vorgestellt. Anschließend erfolgt in Abschnitt \ref{sec:3:qn} eine Darstellung der quantitaven Methode zur Untersuchung des Forschungsgegenstandes und in Abschnitt \ref{sec:3:ql} eine Skizzierung der Analyse der Verbgefüge, die im Korpus vorliegen.

In Abschnitt \ref{sec:4} werden die Ergebnisse aus Abschnitt \ref{sec:3} vorgestellt und anschließend vor zur Beantwortung der Forschungsfrage in Abschnitt \ref{sec:5} diskutiert. Ebenso soll auf die Limitationen der Untersuchung eingegangen werden. Abschnitt \ref{sec:6} fasst die Arbeit zusammen.

